% ==============================================================================
% CAPÍTULO 6 - CONSIDERAÇÕES FINAIS
% ==============================================================================

\chapter{Considerações Finais}
\label{cap:consideracoes}

O estágio realizado na BIOTECNO Indústria e Comércio Ltda. teve como objetivo central o desenvolvimento
de um módulo de geração automatizada de propostas técnicas e comerciais em PDF, com suporte a layouts
intercambiáveis gerenciáveis pelo painel administrativo do Django. O objetivo foi alcançado: ao término
do período, o sistema estava implementado e funcional em ambiente de desenvolvimento, com pipeline de
geração completo da resolução do layout ativo à entrega do PDF compilado pelo LuaLaTeX, pronto
para ser integrado ao ambiente de produção na próxima atualização da aplicação.

%-------------------------------------------------
\section{Conclusões}
%-------------------------------------------------
\label{sec:conclusoes}

A experiência de estágio proporcionou uma imersão em decisões técnicas que raramente se apresentam com
a mesma complexidade e consequência no ambiente acadêmico. A necessidade de integrar dois sistemas com
gramáticas conflitantes, o motor de templates do Django e a sintaxe do LaTeX, exigiu compreensão
profunda de como cada camada funciona internamente, indo além do uso superficial das ferramentas. Da
mesma forma, a decisão de não acoplar o layout ao modelo \texttt{Orcamento} foi fruto de uma análise
de \textit{trade-offs} reais: entre consistência histórica e flexibilidade operacional, optou-se pela
flexibilidade — decisão validada pelo contexto de negócio da empresa.

Do ponto de vista técnico, o estágio consolidou conhecimentos em Django ORM, modelagem relacional,
geração de documentos com LaTeX e integração de processos externos via \texttt{subprocess}, conteúdos
abordados no curso de Ciência da Computação e que, no contexto prático, revelaram nuances e
interdependências não evidentes na teoria. A leitura e compreensão de um sistema de software já em
produção, com múltiplos módulos, convenções próprias e histórico de decisões acumulado, foi por si só
um aprendizado valioso para a formação profissional.

Do ponto de vista comportamental, a dinâmica de trabalho com ciclos curtos de revisão junto ao
supervisor reforçou a importância da comunicação técnica clara e da validação incremental de decisões
de arquitetura. Erros foram identificados e corrigidos precocemente, evitando retrabalho de maior
escala.

Um aspecto relevante que o estágio evidenciou, e que raramente é abordado no ambiente acadêmico,
é a prática de substituição progressiva de código legado por soluções mais modernas e padronizadas.
A migração do módulo de orçamentos do PyLaTeX para a arquitetura de templates Django não foi motivada
por uma falha crítica do sistema anterior, mas pela necessidade de padronização interna e pela
oportunidade de tornar o sistema mais sustentável a longo prazo. Identificar esse tipo de demanda,
propor uma solução arquitetural coerente e executá-la sem quebrar o que já funcionava é uma habilidade
que a formação acadêmica prepara teoricamente, mas que só se desenvolve plenamente na prática.

Por fim, o fato de já atuar na empresa de forma remunerada antes do início do estágio
curricular teve impacto direto na produtividade e na qualidade das entregas. O conhecimento prévio
do sistema Bioweb, das convenções de código adotadas pela equipe e da estrutura do projeto eliminou
a curva de ambientação inicial, permitindo que o desenvolvimento começasse com foco imediato no
problema técnico. Essa familiaridade também facilitou a leitura de módulos existentes como referência, como o uso do \texttt{viewsinspecao.py} para resolver o conflito de delimitadores Django e
\LaTeX{}, acelerando decisões que, em outro contexto, exigiriam investigação mais prolongada.

%-------------------------------------------------
\section{Trabalhos Futuros}
%-------------------------------------------------
\label{sec:trabalhos-futuros}

A continuidade natural do trabalho desenvolvido aponta para os seguintes encaminhamentos:

\begin{itemize}
    \item \textbf{Publicação em produção}: integrar o novo sistema de geração ao ambiente de produção
    da aplicação, substituindo definitivamente a função legada baseada em PyLaTeX, após validação
    completa das variáveis de contexto com dados reais;
    \item \textbf{Múltiplos layouts por tipo de venda}: ampliar o modelo \texttt{LayoutOrcamento} para
    suportar diferentes layouts por categoria de venda (direta, licitação, revenda), permitindo que a
    seleção do template seja automatizada com base nos atributos do orçamento;
    \item \textbf{Interface de pré-visualização}: desenvolver uma funcionalidade de pré-visualização
    do PDF no painel administrativo antes da ativação de um novo layout, reduzindo o risco de
    publicação de templates com erros de formatação.
\end{itemize}